% ============================================================================
% CHAPITRE 3 — SECTION 3.2 : PROTOCOLES EXPÉRIMENTAUX
% Fichier : partie3/protocole_experimental.tex
% ============================================================================

\section{PROTOCOLES EXP\'ERIMENTAUX}
\label{sec:protocoles_experimentaux}
% ============================================================================
\begin{sloppypar}
Cette section sp\'ecifie les six protocoles de mesure destin\'es \`a estimer
empiriquement les param\`etres du cadre CRO, \`a savoir la capacit\'e
d'extraction institutionnelle, l'entropie contextuelle, l'opposabilit\'e,
ainsi que leurs extensions quantiques, en d\'etaillant pour chacun le design,
la proc\'edure, les m\'etriques et les pr\'edictions falsifiables, avant de
consolider le calendrier, les ressources et les crit\`eres de validation de
l'ensemble du programme.
\end{sloppypar}

\begin{boiteArbitre}
L'arbitre ouvre un dossier qui ne ressemble pas aux pr\'ec\'edents, puisqu'il
y trouve des protocoles sans r\'esultats, puisque rien n'a encore \'et\'e mesur\'e.
Il lit la mani\`ere dont on compte mesurer, qui est pr\'ecis\'ement ce qu'il
doit juger.
\end{boiteArbitre}

% -----------------------------------------------------------------------
\subsection{Protocole de mesure de $\Delta$~:
          capacit\'e d'extraction institutionnelle}
\label{subsec:protocole_delta}
% -----------------------------------------------------------------------
\begin{sloppypar}

\subsubsection{Objectif et d\'efinition op\'erationnelle}
\label{subsubsec:objectif_delta}

\noindent\textbf{Objectif.} Mesurer empiriquement la capacit\'e d'extraction
institutionnelle $\Delta$ (en bits), d\'efinie par l'hypoth\`ese d'extraction
born\'ee~\cite{MINKA2026CRO}~:
\begin{equation}
I(V;\Sigma|C) \leq I(M;\Sigma|C) + \Delta
\label{eq:extraction_bornee_op}
\end{equation}

\begin{definition}[Capacit\'e d'extraction institutionnelle $\Delta$]\label{def:capacite_extraction_op}
$\Delta$ est la quantit\'e d'information sur les verdicts qu'un acteur
institutionnel peut extraire \emph{au-del\`a} de l'information mutuelle brute
contenue dans la preuve cryptographique. En pratique, $\Delta$ capture la
connaissance pr\'ealable (pr\'ec\'edents l\'egaux, formation sp\'ecialis\'ee,
proc\'edures institutionnelles, capacit\'e de d\'elib\'eration collective) non
formalis\'ee dans le contexte $C$.
\end{definition}

\noindent\textbf{Estimateur.} Par la
Proposition~\ref{prop:fondement_extraction}, la quantit\'e contr\^ol\'ee par
l'Hypoth\`ese~\ref{ass:extraction_bornee} est l'information conditionnelle
$I(V;\Sigma|M,C)$. L'estimateur retenu est donc
\begin{equation}
\hat{\Delta} = \hat{I}(V;\Sigma \mid M, C)
\label{eq:estimateur_delta}
\end{equation}
o\`u $\hat{I}$ est un estimateur de l'information mutuelle conditionnelle
(plug-in avec correction de Miller--Madow, ou estimateur de type NSB pour
les petits \'echantillons), calcul\'e sur les triplets $(v_i, \sigma_i, m_i)$
recueillis \`a contexte $c$ fix\'e. Le message $m_i$ est connu de
l'exp\'erimentateur (v\'erit\'e terrain) mais non des participants. $\hat{\Delta}$ estime
la plus petite capacit\'e compatible avec l'Hypoth\`ese~\ref{ass:extraction_bornee}
sous la condition~\eqref{eq:extraction_bornee_forte}. L'estimateur est biais\'e
positivement en petit \'echantillon~; l'intervalle de confiance est obtenu par
bootstrap non param\'etrique \emph{sur les participants}.

\begin{remark}[Formulation initiale de l'estimateur]
La formulation initiale de ce protocole proposait
$\hat{\Delta} = I_{\mathrm{emp}}(V;\Sigma|C) - I_{\mathrm{vrai}}(M;\Sigma|C)$.
Cet estimateur prend des valeurs n\'egatives pour un protocole fiable
($I(M;\Sigma|C) \geq \Relrec\cdot\lambda - 1$ par Fano, alors que
$I(V;\Sigma|C) \leq \log_2|\mathcal{V}_\lambda|$) et ne fournit qu'une borne
inf\'erieure sur $\Delta$. La formulation retenue
est~\eqref{eq:estimateur_delta}.
\end{remark}

\subsubsection{Design exp\'erimental}
\label{subsubsec:design_delta}

\noindent\textbf{Participants.} $N \geq 100$ professionnels des institutions
cibles, stratifi\'es par niveau d'expertise et par type
d'institution~\cite{MINKA2026CRO} selon la r\'epartition du
tableau~\ref{tab:stratification_delta}.

\begin{table}[htbp]
\centering
\caption{Stratification de l'\'echantillon pour la mesure de $\Delta$}
\label{tab:stratification_delta}
\begin{tabular}{@{}llcc@{}}
\toprule
\textbf{Cat\'egorie} & \textbf{Sous-groupe} & \textbf{$n$} & \textbf{$\Delta$ pr\'edit} \\
\midrule
Juges & G\'en\'eralistes (juridictions civiles) & 20 & $\approx 25$ bits \\
      & Sp\'ecialis\'es (contrats, PKI, num\'erique) & 20 & $\approx 40$ bits \\
\midrule
R\'egulateurs & Juniors (< 5 ans d'exp\'erience) & 20 & $\approx 35$ bits \\
financiers  & Seniors ($\geq 5$ ans, superviseurs) & 20 & $\approx 65$ bits \\
\midrule
Auditeurs & Cour des comptes (secteur public) & 20 & $\approx 45$ bits \\
          & Cabinets priv\'es (Big 4 et \'equivalents) & 20 & $\approx 55$ bits \\
\midrule
\textbf{Total} & & \textbf{120} & \\
\bottomrule
\end{tabular}
\end{table}

\noindent\textbf{Mat\'eriel exp\'erimental.}
\begin{itemize}
\item 5 types de preuves cryptographiques (ECDSA, Helios avec ElGamal et ZK,
      Groth16 zk-SNARK, RSA-PSS et ElGamal classique)
\item 10 contextes standardis\'es par type de preuve (contractuels,
      \'electoraux, financiers), soit 50 paires contexte--preuve
\item Questionnaires valid\'es de 30 items par contexte
\end{itemize}

\subsubsection{Proc\'edure}
\label{subsubsec:procedure_delta}

La session exp\'erimentale pour chaque participant se d\'eroule en cinq
phases.

\noindent\textbf{Phase 0. Formation initiale (2 h).}
Introduction aux concepts cryptographiques de base sans r\'ev\'eler les
m\'ecanismes internes, afin d'homog\'en\'eiser le vocabulaire minimum.

\noindent\textbf{Phase 1. Session exp\'erimentale principale (4 h).}
Pour chacun des 10 contextes assign\'es (rotation \'equilibr\'ee)~:
\begin{enumerate}
\item[(a)] Pr\'esentation du contexte $c \in \mathcal{C}$
\item[(b)] Pr\'esentation de la preuve cryptographique $\sigma$
\item[(c)] Compl\'etion du questionnaire d'extraction de verdict
\item[(d)] Mesure du temps de r\'eponse et notation de confiance subjective
\end{enumerate}

\noindent\textbf{Phase 2. V\'erification test-retest (1 h).}
Re-pr\'esentation de 3 contextes d\'ej\`a vus apr\`es un d\'elai de 2 semaines.

\noindent\textbf{Phase 3. Calcul du gold standard.}
Deux cryptographes experts ind\'ependants rendent les verdicts \og corrects\fg{}
(accord inter-juge requis $\geq 90\%$).

\noindent\textbf{Phase 4. D\'ebriefing et restitution (30 min).}

\subsubsection{M\'etriques et estimateurs}
\label{subsubsec:metriques_delta}

Le protocole standardis\'e de~\cite{MINKA2026CRO} calcule $\hat\Delta$ selon
l'Algorithme~\ref{alg:estimation_delta}.

\begin{algorithm}[H]
\caption{Estimation empirique de $\Delta$}
\label{alg:estimation_delta}
\begin{algorithmic}[1]
\State \textbf{Entr\'ees~:} Institution $I$, types de preuves
       $\{\sigma_1, \ldots, \sigma_k\}$, contexte $C$
\State \textbf{Sortie~:} $\hat\Delta$, $\hat\alpha$
\State Pr\'esenter $\sigma_i$ avec $C$ aux acteurs institutionnels
\State Enregistrer verdicts $V$, scores de confiance, temps de d\'ecision
\State Associer \`a chaque r\'eponse le message $m$ connu de l'exp\'erimentateur
\State Tester l'admissibilit\'e $V \perp M \mid (\Sigma, C)$ (test de permutation)
\State Estimer $\hat\Delta = \hat I(V;\Sigma \mid M, C)$ (\'equation~\eqref{eq:estimateur_delta})
\State Estimer $\hat\alpha = \log_2|\mathcal{V}| / H(C)$
\State Valider par fiabilit\'e test-retest et coh\'erence inter-\'evaluateurs
\end{algorithmic}
\end{algorithm}

L'algorithme transpose sur le terrain la d\'efinition th\'eorique de la
capacit\'e d'extraction institutionnelle en encha\^inant quatre phases.

La phase d'\textbf{entr\'ees} (lignes~1--2) fixe le cadre de la mesure.
Elle re\c{c}oit une institution cible $I$, un ensemble de types de preuves
cryptographiques $\{\sigma_1, \ldots, \sigma_k\}$ et un contexte $C$, puis
d\'eclare comme sorties l'estimateur $\hat\Delta$ de la capacit\'e
d'extraction et l'estimateur $\hat\alpha$ de la granularit\'e des verdicts.
Aucun calcul n'intervient \`a ce stade, qui se borne \`a nommer les grandeurs
mesur\'ees. S\'eparer entr\'ees et sorties importe n\'eanmoins, car $\Delta$
n'est jamais observ\'e directement et se d\'eduit d'un contraste entre
deux quantit\'es informationnelles.

La phase de \textbf{collecte} (lignes~3--4) produit les donn\'ees brutes.
Elle pr\'esente chaque preuve $\sigma_i$ accompagn\'ee de son contexte $C$
aux acteurs institutionnels, puis enregistre les verdicts $V$ rendus, les
scores de confiance subjectifs et les temps de d\'ecision, trois s\'eries qui
doivent \^etre recueillies ensemble. Les verdicts seuls
suffiraient \`a calculer une distribution marginale, mais les scores de
confiance et les temps de d\'ecision permettent de d\'etecter les r\'eponses
prises au hasard, qui contamineraient l'estimation de l'information
mutuelle si elles \'etaient trait\'ees comme des jugements inform\'es.

La phase d'\textbf{appariement et de contr\^ole} (lignes~5--6) associe
chaque verdict au message effectivement engag\'e, que l'exp\'erimentateur
conna\^it et que les participants ignorent. Elle teste ensuite la condition
d'admissibilit\'e de la D\'efinition~\ref{def:interp_admissible}. Si $V$
d\'epend encore de $M$ une fois $(\Sigma,C)$ connus, les participants ont
acc\`es \`a une information hors preuve de sorte que l'\'etude sort du cadre.

La phase d'\textbf{estimation} (lignes~7--8) produit les deux param\`etres.
L'estimateur $\hat\Delta = \hat I(V;\Sigma|M,C)$ mesure ce que les acteurs
tirent de la preuve \emph{au-del\`a} de ce qu'elle dit du message, qu'il s'agisse
de la mise en forme, des m\'etadonn\'ees ou des indices que l'expertise institutionnelle
(pr\'ec\'edents l\'egaux, expertise de domaine, proc\'edures internes) rend
lisibles. Ce terme co\"{\i}ncide avec celui que la
Proposition~\ref{prop:fondement_extraction} isole dans la d\'ecomposition de
$I(V;\Sigma|C)$. L'estimateur $\hat\alpha$ est
calcul\'e ind\'ependamment comme le rapport
$\log_2|\mathcal{V}| / H(C)$ et mesure la finesse de l'espace de verdicts
relativement \`a la richesse contextuelle, conform\'ement \`a la
d\'efinition~\ref{def:alpha_granularite}. Les deux estimateurs sont donc
acquis par des voies distinctes, ce qui permet de les valider
s\'epar\'ement.

La phase de \textbf{validation} (ligne~9) contr\^ole la fiabilit\'e de la
mesure. Elle combine la stabilit\'e test-retest, qui
v\'erifie que les m\^emes acteurs produisent des r\'eponses coh\'erentes
lorsqu'ils sont r\'einterrog\'es apr\`es un d\'elai, avec la coh\'erence
inter-\'evaluateurs, qui v\'erifie que des acteurs du m\^eme profil
produisent des estimations convergentes. Sans cette \'etape, une valeur
\'elev\'ee de $\hat\Delta$ pourrait refl\'eter une simple variabilit\'e
al\'eatoire au lieu d'une capacit\'e d'extraction r\'eelle.

L'algorithme fournit ainsi une estimation de $\Delta$ dont la stabilit\'e est
contr\^ol\'ee, sans pr\'etendre \`a une valeur exacte. Le cadre CRO n'en demande
pas davantage, puisque $\Delta$ y intervient
comme constante du th\'eor\`eme et doit \^etre mesur\'e avec une incertitude
born\'ee pour que les pr\'edictions aient un sens op\'erationnel.

\subsubsection{Pr\'edictions v\'erifiables}
\label{subsubsec:predictions_delta}

Les pr\'edictions suivantes sont directement
falsifiables~\cite{MINKA2026CRO,MINKA2026SOTP}~:

\begin{enumerate}
\item $\hat\Delta_{\text{juge g\'en\'eral}} \approx 25$ bits
\item $\hat\Delta_{\text{juge PKI}} \approx 40$ bits
\item $\hat\Delta_{\text{r\'egulateur senior}} \approx 65$ bits
\item $\hat\Delta_{\text{ECDSA}} > \hat\Delta_{\text{Helios}}$ pour tout sous-groupe
\item Facteur d'amplification institutionnelle $1{,}8$ \`a $3{,}7$ fois par
      rapport \`a la capacit\'e individuelle Miller ($\approx 20$ bits)
\end{enumerate}
\end{sloppypar}

\begin{boiteArbitre}
L'arbitre note un \'el\'ement nouveau, \`a savoir des chiffres qui fixent des cibles,
avec vingt-cinq bits pour un juge g\'en\'eraliste et soixante-cinq pour un
r\'egulateur senior. Il y voit des paris, qui peuvent donc \^etre perdus.
\end{boiteArbitre}

% -----------------------------------------------------------------------
\subsection{Protocole de mesure de $H(\mathcal{C})$~:
          entropie contextuelle}
\label{subsec:protocole_HC}
% -----------------------------------------------------------------------
\begin{sloppypar}

\subsubsection{Objectif et enjeu}
\label{subsubsec:objectif_HC}

\noindent\textbf{Objectif.} Quantifier empiriquement l'entropie contextuelle
$H(\mathcal{C})$, c'est-\`a-dire la richesse informationnelle du contexte
institutionnel dans lequel s'op\`ere l'interpr\'etation des preuves
cryptographiques.

\noindent\textbf{Enjeu.} $H(\mathcal{C})$ est le param\`etre le plus difficile
\`a mesurer directement, car il caract\'erise la \emph{distribution}
des contextes possibles et non un contexte particulier. Cette distinction
d'\'echelle est explicit\'ee \`a la
section~\ref{sec:methodologie_calcul}. Les protocoles ci-dessous mesurent
l'entropie de la distribution des contextes d'un domaine institutionnel,
tandis que les d\'erivations de la section~\ref{sec:derivation_opposabilite}
mobilisent l'entropie du contexte d'esp\`ece. Les deux grandeurs portent le
m\^eme symbole et ne prennent pas les m\^emes valeurs.

\subsubsection{Approche 1~: d\'ecomposition structurelle}
\label{subsubsec:approche_structurelle}

\noindent\textbf{Principe.} Identifier les dimensions ind\'ependantes du
contexte institutionnel et sommer leurs entropies en supposant
l'ind\'ependance approximative~\cite{MINKA2026CLO}.

\noindent\textbf{Application au contexte contractuel.}
\begin{equation}
H(C) = H(C_{\text{parties}}) + H(C_{\text{clauses}}) + H(C_{\text{juridiction}}) + H(C_{\text{historique}})
\label{eq:decomposition_HC}
\end{equation}

\begin{table}[htbp]
\centering
\caption{Estimations structurelles de $H(\mathcal{C})$ par domaine institutionnel}
\label{tab:estimations_structurelles_HC}
\begin{tabular}{@{}llc@{}}
\toprule
\textbf{Domaine} & \textbf{Composantes dominantes} & $H(\mathcal{C})$ estim\'e \\
\midrule
Contentieux \'electoral & Registre, r\`egles, historique & 30--50 bits \\
Conformit\'e financi\`ere & Transactions, r\'eglementations, audit & 40--70 bits \\
Signature contractuelle & Parties, clauses, juridiction & 20--40 bits \\
\bottomrule
\end{tabular}
\end{table}

\subsubsection{Approche 2~: mod\'elisation probabiliste par MDS}
\label{subsubsec:approche_MDS}

\noindent\textbf{Design.} Pr\'esenter un ensemble de $n_{\text{ctx}} = 100$
contextes vari\'es \`a un panel de 5--10 experts en leur demandant une
\'evaluation de similarit\'e par paires.

\noindent\textbf{Proc\'edure d'analyse.}
\begin{enumerate}
\item Construction de la matrice de dissimilarit\'e $D \in \mathbb{R}^{100 \times 100}$
\item R\'eduction par MDS (multi-dimensional scaling)
\item Calcul du volume effectif $V_{\mathrm{eff}} = \det(\hat\Sigma)^{1/2}$
\item Estimation $H(\mathcal{C}) \approx \log_2 V_{\mathrm{eff}}$
\end{enumerate}

\subsubsection{Approche 3~: consensus d'experts par m\'ethode Delphi}
\label{subsubsec:approche_delphi}

\noindent\textbf{Principe.} La m\'ethode Delphi est une technique
d'\'elicitation it\'erative permettant d'obtenir un consensus
quantitatif~\cite{MINKA2026CRO}.

\noindent\textbf{Proc\'edure en quatre tours.} 20 experts estiment
$H(\mathcal{C})$ ind\'ependamment, puis r\'evisent apr\`es retour
anonymis\'e. Le processus s'arr\^ete lorsque l'IQR devient inf\'erieur \`a $5$ bits.
\end{sloppypar}

% -----------------------------------------------------------------------
\subsection{Protocole de mesure de $\Opp$~:
          opposabilit\'e sur syst\`emes d\'eploy\'es}
\label{subsec:protocole_opp}
% -----------------------------------------------------------------------
\begin{sloppypar}

\subsubsection{Objectif et syst\`eme cible}
\label{subsubsec:objectif_opp}

\noindent\textbf{Objectif.} Mesurer l'opposabilit\'e $\Opp$ empiriquement sur
un syst\`eme cryptographique r\'eellement d\'eploy\'e puis comparer la valeur
mesur\'ee \`a la pr\'ediction du mod\`ele SOTP~\cite{MINKA2026SOTP}.

\noindent\textbf{Syst\`eme cible.} Syst\`eme de vote Helios d\'eploy\'e pour
une \'election de d\'el\'egu\'es \'etudiants \`a l'ENSP (\'Ecole Nationale
Sup\'erieure Polytechnique, Universit\'e de Yaound\'e I).

\subsubsection{Design et participants}
\label{subsubsec:design_opp}

\noindent\textbf{Participants.} Commission \'electorale de $n = 10$ trustees
(\'etudiants doctorants et enseignants-chercheurs).

\noindent\textbf{Taille de l'\'election.} $\approx 500$ \'etudiants votants,
g\'en\'erant $\approx 500$ preuves cryptographiques.

\subsubsection{Proc\'edure}
\label{subsubsec:procedure_opp}

\noindent\textbf{Phase 1. Formation trustees (4 h).}
Formation structur\'ee sur Helios et les preuves z\'ero-connaissance.

\noindent\textbf{Phase 2. \'Election r\'eelle.}
D\'eploiement standard de Helios.

\noindent\textbf{Phase 3. Session d'audit post-\'election (3 h).}
Pour un \'echantillon al\'eatoire de 50 preuves (10\%)~:
\begin{enumerate}
\item Chaque trustee examine ind\'ependamment la preuve
\item Rend un verdict $v \in \mathcal{V}$
\item Note sa confiance sur 10
\end{enumerate}

\subsubsection{Calcul de l'opposabilit\'e empirique}
\label{subsubsec:calcul_opp_empirique}

L'opposabilit\'e est calcul\'ee selon~\cite{MINKA2026CLO}~:
\begin{equation}
\Opp_{\mathrm{emp}} = \frac{I_{\mathrm{emp}}(V;\Sigma|C)}{H_{\mathrm{emp}}(V|C)}
\label{eq:opp_empirique}
\end{equation}

\subsubsection{Validation crois\'ee avec la pr\'ediction SOTP}
\label{subsubsec:validation_croisee_sotp}

\noindent\textbf{Valeur de r\'ef\'erence.} Pour Helios, la calibration
illustrative de l'article CLO~\cite{MINKA2026CLO} donne, \`a l'\'echelle
op\'erationnelle (fiabilit\'e de v\'erification)~:
\begin{equation}
\Conf \cdot \Relver \cdot \Opp \approx 0{,}437
\label{eq:prediction_sotp_opp}
\end{equation}
La valeur cible pour l'opposabilit\'e seule est
$\Opp_{\mathrm{pr\'edit}} = 0{,}50$~; sa d\'erivation compl\`ete (espace de
verdicts \`a quatre d\'ecisions, signal expert binaire d'un bit) est
\'etablie \`a la section~\ref{sec:derivation_helios}. Le
pr\'esent protocole met \`a l'\'epreuve l'hypoth\`ese d'uniformit\'e des verdicts qui
sous-tend cette valeur. Il mesure la distribution empirique $H_{\mathrm{emp}}(V|C)$
au lieu de la postuler.

\noindent\textbf{Crit\`ere de validation.}
\begin{itemize}
\item Accord si $|\Opp_{\mathrm{emp}} - \Opp_{\mathrm{pr\'edit}}| \leq 0{,}05$
\item D\'esaccord (la valeur $0{,}50$ et l'hypoth\`ese d'uniformit\'e qui la fonde
      sont r\'efut\'ees) si $\Opp_{\mathrm{pr\'edit}}$ tombe hors de l'IC~95\,\% bootstrap
      de $\Opp_{\mathrm{emp}}$
\end{itemize}
\end{sloppypar}

% -----------------------------------------------------------------------
\subsection{Protocole standardis\'e d'estimation conjointe de $\Delta$ et $\alpha$}
\label{subsec:protocole_conjoint}
% -----------------------------------------------------------------------
\begin{sloppypar}
Pour une validation prospective compl\`ete du cadre CRO, les
param\`etres $(\Delta, \alpha, \beta)$ doivent \^etre mesur\'es
\emph{ind\'ependamment} et \emph{conjointement} avant que les pr\'edictions
CRO ne soient calcul\'ees.

\begin{algorithm}[H]
\caption{Protocole standardis\'e d'estimation conjointe de
         $(\hat\Delta, \hat\alpha)$~\cite{MINKA2026CRO}}
\label{alg:estimation_conjointe}
\begin{algorithmic}[1]
\State \textbf{Entr\'ees~:} Institution $I$, types de preuves, contexte $C$
\State \textbf{Sorties~:} $\hat\Delta$ (bits), $\hat\alpha$ (sans unit\'e)
\Require Accord institutionnel, \'ethique IRB, $n \geq 30$ participants
\State \textbf{\'Etape 1~:} Collecte de donn\'ees (verdicts, confiance, temps)
\State \textbf{\'Etape 2~:} Appariement des verdicts au message $m$ (v\'erit\'e terrain)
\State \textbf{\'Etape 3~:} Test d'admissibilit\'e $V \perp M \mid (\Sigma,C)$
\State \textbf{\'Etape 4~:} $\hat\Delta \leftarrow \hat I(V;\Sigma \mid M,C)$
\State \textbf{\'Etape 5~:} $\hat\alpha \leftarrow \log_2|\mathcal{V}| / H(C)$
\State \textbf{\'Etape 6~:} Validation (ICC $\geq 0{,}8$, $\kappa \geq 0{,}7$)
\end{algorithmic}
\end{algorithm}

L'algorithme formalise en six \'etapes le protocole de r\'ef\'erence pour
l'estimation conjointe des param\`etres institutionnels $(\Delta, \alpha)$,
en int\'egrant explicitement les pr\'ealables \'ethiques et
m\'ethodologiques que la version op\'erationnelle de la
section~\ref{subsec:protocole_delta} suppose implicites.

La phase de \textbf{param\'etrage} (lignes~1--3) fixe le cadre de
l'exp\'erience. Elle d\'eclare l'institution cible $I$, les types de
preuves \`a pr\'esenter et le contexte $C$, puis nomme les deux sorties,
$\hat\Delta$ en bits et $\hat\alpha$ sans unit\'e. Elle pose surtout trois
\emph{pr\'erequis} obligatoires, \`a savoir l'accord institutionnel pr\'ealable,
l'approbation d'un comit\'e d'\'ethique de la recherche et un effectif
minimal de $n \geq 30$ participants. Leur port\'ee d\'epasse la simple formalit\'e
proc\'edurale, puisqu'ils conditionnent la validit\'e statistique de
l'estimation. En de\c{c}\`a de trente participants, les estimateurs
d'information mutuelle deviennent instables~; l'approbation \'ethique
autorise quant \`a elle la collecte de donn\'ees sur des acteurs humains
identifiables. La ligne~3, marqu\'ee \texttt{Require}, signale que ces
conditions doivent \^etre satisfaites \emph{avant} toute ex\'ecution, car elles
sont des conditions d'entr\'ee de l'algorithme, distinctes de ses \'etapes.

L'\textbf{\'etape~1} (ligne~4) collecte les donn\'ees brutes. Elle
enregistre, pour chaque participant et chaque paire contexte-preuve, les
verdicts rendus, les scores de confiance subjectifs et les
temps de d\'ecision. Ces trois s\'eries sont reprises \`a l'identique de
l'algorithme~\ref{alg:estimation_delta}, dont le protocole standardis\'e
conserve la nature en se bornant \`a en fixer la s\'equence et les crit\`eres
de validit\'e. Les scores de confiance et les temps servent de variables
de contr\^ole pour d\'etecter les r\'eponses non inform\'ees.

Les \textbf{\'etapes~2 et 3} (lignes~5--6) apparient chaque verdict au
message engag\'e, connu du seul exp\'erimentateur, puis v\'erifient que les
participants n'ont pas acc\`es \`a $M$ hors de la preuve (condition
d'admissibilit\'e). Cette v\'erification conditionne l'interpr\'etation de
l'\'etape suivante.

Les \textbf{\'etapes~4 et 5} (lignes~7--8) produisent les deux
estimateurs. L'\'etape~4 applique l'estimateur~\eqref{eq:estimateur_delta}
$\hat\Delta \leftarrow \hat I(V;\Sigma|M,C)$.
L'\'etape~5 calcule $\hat\alpha$ par la formule
$\log_2|\mathcal{V}| / H(C)$, ind\'ependamment de $\hat\Delta$. Les deux calculs sont d\'elib\'er\'ement s\'epar\'es afin de valider
chaque estimateur selon ses propres crit\`eres, sans que l'erreur sur
l'un ne contamine l'autre~; un $\hat\Delta$ biais\'e n'entra\^ine donc pas
m\'ecaniquement un $\hat\alpha$ biais\'e.

L'\textbf{\'etape~6} (ligne~9) valide les estimations par deux
statistiques. Le coefficient de corr\'elation intra-classe
$\mathrm{ICC} \geq 0{,}8$ mesure la coh\'erence test-retest tandis que le
coefficient $\kappa$ de Cohen $\geq 0{,}7$ mesure l'accord
inter-\'evaluateurs au-del\`a du hasard. Ces deux seuils sont plus
exigeants que ceux de l'algorithme~\ref{alg:estimation_delta}, qui se
contentait d'une validation qualitative. Le protocole standardis\'e fixe
donc des crit\`eres num\'eriques pr\'ecis qui le rendent reproductible
par des \'equipes ind\'ependantes. Cette reproductibilit\'e justifie le qualificatif
de \emph{standardis\'e}, puisque deux
laboratoires appliquant les m\^emes \'etapes aux m\^emes entr\'ees doivent
obtenir des estimations compatibles.

Par rapport \`a l'algorithme~\ref{alg:estimation_delta}, cette version
explicite les pr\'erequis \'ethiques et statistiques, num\'erote les \'etapes
pour permettre de les citer individuellement et fixe des seuils quantitatifs
de validation. Ces ajouts font de la proc\'edure de mesure un protocole
standardis\'e, r\'eutilisable au-del\`a du cas d'esp\`ece.
\end{sloppypar}

% -----------------------------------------------------------------------
\subsection{Protocole de mesure de $\Delta_q$~:
          capacit\'e d'extraction quantique}
\label{subsec:protocole_delta_q}
% -----------------------------------------------------------------------
\begin{sloppypar}

\subsubsection{Objectif et d\'efinition op\'erationnelle}
\label{subsubsec:objectif_delta_q}

\noindent\textbf{Objectif.} Mesurer empiriquement la capacit\'e d'extraction
institutionnelle quantique $\Delta_q$ (en bits), d\'efinie
par~\cite{MINKA2026PQCRO}~:
\begin{equation}
I(V^q; \Sigma|C)_\rho \leq I(M;\Sigma|C)_\rho + \Delta_q
\label{eq:extraction_quantique_op}
\end{equation}

\begin{definition}[Capacit\'e d'extraction quantique $\Delta_q$]\label{def:capacite_extraction_quantique_op}
$\Delta_q$ est la quantit\'e d'information sur les verdicts qu'un acteur
institutionnel peut extraire dans un environnement post-quantique,
au-del\`a de l'information accessible de Holevo. Elle se d\'ecompose
comme $\Delta_q = \Delta + \delta_{\mathrm{enh}} - \delta_{\mathrm{degr}}$,
o\`u $\delta_{\mathrm{enh}}$ repr\'esente l'am\'elioration due aux
ressources quantiques et $\delta_{\mathrm{degr}}$ la d\'egradation due au
bruit quantique.
\end{definition}

\noindent\textbf{Estimateur.}
\begin{equation}
\hat{\Delta}_q = \hat I(V^q;\Sigma \mid M, C)
\label{eq:estimateur_delta_q}
\end{equation}
calcul\'e comme~\eqref{eq:estimateur_delta} \`a partir des r\'eponses des
participants, le message $m$ \'etant connu de l'exp\'erimentateur. Pour des
signatures post-quantiques qui sont des cha\^{\i}nes de bits classiques, cet
estimateur est identique \`a l'estimateur classique~; il ne diff\`ere que dans le
sc\'enario hypoth\'etique d'un support ou d'un v\'erificateur quantique.

\subsubsection{Design exp\'erimental}
\label{subsubsec:design_delta_q}

\noindent\textbf{Participants.} $N \geq 80$ professionnels form\'es aux
syst\`emes post-quantiques, stratifi\'es par type d'institution selon la
r\'epartition du tableau~\ref{tab:stratification_delta_q}.

\begin{table}[htbp]
\centering
\caption{Stratification de l'\'echantillon pour la mesure de $\Delta_q$}
\label{tab:stratification_delta_q}
\begin{tabular}{@{}llcc@{}}
\toprule
\textbf{Cat\'egorie} & \textbf{Sous-groupe} & \textbf{$n$} & \textbf{$\Delta_q$ pr\'edit} \\
\midrule
R\'egulateurs PQC & Juniors (formation initiale) & 20 & $\approx 28$ bits \\
                  & Seniors (formation avanc\'ee) & 20 & $\approx 45$ bits \\
\midrule
Auditeurs crypto & Classiques (transition PQC) & 20 & $\approx 35$ bits \\
                 & Experts PQC certifi\'es & 20 & $\approx 67$ bits \\
\midrule
\textbf{Total} & & \textbf{80} & \\
\bottomrule
\end{tabular}
\end{table}

\noindent\textbf{Mat\'eriel exp\'erimental.}
\begin{itemize}
\item 2 sch\'emas de signature post-quantiques, ML-DSA-65 (FIPS~204,
      ex-Dilithium-3) et SLH-DSA (FIPS~205, ex-SPHINCS+)~; ML-KEM (FIPS~203,
      ex-Kyber) est un m\'ecanisme d'encapsulation de cl\'e et ne produit pas de
      preuve, il est donc exclu
\item Simulateur d'environnement quantique avec bruit de d\'epolarisation
      param\'etrable ($p \in [0{,}001;\,0{,}1]$)
\item Interface de v\'erification post-quantique avec m\'etriques de performance
\item Questionnaires adapt\'es au contexte NIST PQC (FIPS 203/204/205)
\end{itemize}

\subsubsection{Proc\'edure}
\label{subsubsec:procedure_delta_q}

La session exp\'erimentale se d\'eroule en six phases.

\noindent\textbf{Phase 0. Formation PQC (3 h).}
Introduction aux standards NIST PQC portant sur la structure des cl\'es, le processus de
signature et les diff\'erences avec RSA/ECDSA. Formation sur les implications
institutionnelles de la transition post-quantique.

\noindent\textbf{Phase 1. Session exp\'erimentale principale (4 h).}
Pour chaque type de preuve post-quantique~:
\begin{enumerate}
\item[(a)] Pr\'esentation du contexte (transaction, document, conformit\'e)
\item[(b)] Pr\'esentation de la signature ML-DSA ou SLH-DSA
\item[(c)] Extraction de verdict avec et sans aide \`a l'interpr\'etation
\item[(d)] Mesure des temps de d\'ecision et scores de confiance
\end{enumerate}

\noindent\textbf{Phase 2. Simulation quantique (2 h).}
Exposition \`a des sc\'enarios simul\'es avec bruit quantique variable~:
\begin{itemize}
\item Condition A, sans bruit (r\'ef\'erence classique)
\item Condition B, bruit faible ($p = 0{,}001$, typique des syst\`emes corrig\'es)
\item Condition C, bruit mod\'er\'e ($p = 0{,}01$, typique NISQ)
\item Condition D, bruit \'elev\'e ($p = 0{,}1$, test de robustesse)
\end{itemize}

\noindent\textbf{Phase 3. Test-retest (1 h).}
Re-pr\'esentation de sc\'enarios apr\`es 2 semaines.

\noindent\textbf{Phase 4. Gold standard PQC.}
Experts en cryptographie post-quantique \'etablissent les verdicts corrects.

\noindent\textbf{Phase 5. D\'ebriefing (30 min).}

\subsubsection{M\'etriques et estimateurs}
\label{subsubsec:metriques_delta_q}

Le protocole calcule $\hat\Delta_q$ et $\hat\eta_q$ selon
l'Algorithme~\ref{alg:estimation_delta_q}.

\begin{algorithm}[H]
\caption{Estimation empirique de $\Delta_q$}
\label{alg:estimation_delta_q}
\begin{algorithmic}[1]
\State \textbf{Entr\'ees~:} Institution $I$, preuves PQC $\{\sigma_1, \sigma_2\}$,
       param\`etre de bruit $p$
\State \textbf{Sortie~:} $\hat\Delta_q$, $\hat\eta_q$
\State Pr\'esenter $\sigma_i$ avec contexte $C$ dans environnement simul\'e
\State Enregistrer verdicts $V^q$, confiance, temps, erreurs d'interpr\'etation
\State Apparier chaque verdict au message $m$ (v\'erit\'e terrain)
\State Tester l'admissibilit\'e $V^q \perp M \mid (\Sigma, C)$
\State Estimer $\hat\Delta_q = \hat I(V^q;\Sigma \mid M, C)$
\State Estimer $\hat\eta_q$ par r\'egression sur les niveaux de bruit
\State Valider par test-retest et accord inter-\'evaluateurs
\end{algorithmic}
\end{algorithm}

L'algorithme transpose au r\'egime post-quantique la proc\'edure
d'estimation de la capacit\'e d'extraction institutionnelle, en y
int\'egrant deux diff\'erences structurelles. L'environnement de
pr\'esentation y est simul\'e avec un param\`etre de bruit quantique
contr\^ol\'e~; la sortie comprend aussi, en plus de $\hat\Delta_q$,
l'estimateur $\hat\eta_q$ du coefficient de perte
d'interpr\'etabilit\'e.

La phase d'\textbf{entr\'ees} (lignes~1--2) fixe le cadre. Elle re\c{c}oit
l'institution cible $I$, un ensemble de deux preuves post-quantiques
$\{\sigma_1, \sigma_2\}$ correspondant aux deux standards de signature NIST
PQC retenus ainsi qu'un param\`etre de bruit $p$ qui contr\^ole l'intensit\'e
de la d\'egradation quantique simul\'ee. Elle d\'eclare comme sorties
l'estimateur $\hat\Delta_q$ de la capacit\'e d'extraction quantique et
l'estimateur $\hat\eta_q$ du coefficient de perte s\'emantique. Cette double
sortie distingue l'algorithme de son homologue classique, puisque
l'algorithme~\ref{alg:estimation_delta} ne mesurait qu'une capacit\'e alors que la
version quantique mesure \'egalement la d\'egradation introduite par le
canal.

La phase de \textbf{collecte} (lignes~3--4) produit les donn\'ees brutes
dans un environnement simul\'e. Chaque preuve $\sigma_i$ est pr\'esent\'ee
avec son contexte $C$, mais contrairement \`a la version classique,
l'environnement de pr\'esentation est param\'etr\'e par $p$, de sorte que les
acteurs institutionnels voient les preuves \`a travers un canal quantique
dont le bruit est contr\^ol\'e. Aux verdicts $V^q$, aux scores de confiance et
aux temps de d\'ecision s'ajoute une quatri\`eme s\'erie, celle des erreurs
d'interpr\'etation. Elle permet de saisir ce que les acteurs extraient
\emph{incorrectement} sous l'effet du bruit, faute de quoi la d\'egradation
s\'emantique resterait invisible dans les donn\'ees.

La phase d'\textbf{appariement et de contr\^ole} (lignes~5--6) reprend la
structure classique, puisque chaque verdict est associ\'e au message engag\'e
avant que l'admissibilit\'e de l'interpr\`ete soit test\'ee. Dans le sc\'enario d'un
support quantique, les quantit\'es sont index\'ees par l'\'etat $\rho$~;
pour des signatures qui sont des cha\^{\i}nes de bits, elles co\"{\i}ncident
avec leurs analogues classiques.

La phase d'\textbf{estimation} (lignes~7--8) produit les deux param\`etres
par des voies distinctes. L'estimateur $\hat\Delta_q = \hat I(V^q;\Sigma|M,C)$ est le
m\^eme que dans le r\'egime classique et quantifie ce que les acteurs
institutionnels tirent de la preuve au-del\`a de ce qu'elle dit du message. L'estimateur $\hat\eta_q$ provient en revanche
d'une r\'egression sur les niveaux de bruit. En r\'ep\'etant
la mesure pour plusieurs valeurs de $p$, on observe comment
l'information extraite d\'ecro\^it avec le bruit~; la pente de cette
d\'ecroissance fournit $\hat\eta_q$. Cette s\'eparation des deux
estimations s'impose, car $\Delta_q$ et $\eta_q$ ont des
interpr\'etations distinctes (l'une mesure une capacit\'e, l'autre une
perte) et doivent \^etre estim\'es ind\'ependamment pour que chacun reste
interpr\'etable dans la borne post-quantique.

La phase de \textbf{validation} (ligne~9) reprend les crit\`eres de
stabilit\'e test-retest et de coh\'erence inter-\'evaluateurs de la
version classique. Ces crit\`eres sont n\'eanmoins plus difficiles \`a
satisfaire en r\'egime quantique, car le bruit introduit une variabilit\'e
suppl\'ementaire dans les r\'eponses. Une faible stabilit\'e test-retest
peut alors refl\'eter soit une incoh\'erence des acteurs, soit une
instabilit\'e intrins\`eque du canal quantique~; distinguer les deux
n\'ecessite de r\'ep\'eter la mesure \`a bruit constant avant de faire
varier $p$.

Par rapport \`a son homologue classique, l'algorithme op\`ere dans un
environnement simul\'e param\'etr\'e, enregistre les erreurs d'interpr\'etation,
calcule des informations mutuelles index\'ees par $\rho$ et produit par
r\'egression un second estimateur $\hat\eta_q$. Ces diff\'erences pr\'eservent
la logique du contraste informationnel, qu'elles \'etendent au r\'egime
quantique en y ajoutant la mesure explicite de la d\'egradation
s\'emantique.

\subsubsection{Pr\'edictions v\'erifiables pour $\Delta_q$}
\label{subsubsec:predictions_delta_q}

\begin{enumerate}
\item $\hat\Delta_q < \hat\Delta$ pour les acteurs non form\'es (p\'enalit\'e d'adaptation)
\item $\hat\Delta_q \to \hat\Delta$ pour les experts PQC (adaptation r\'eussie)
\item $\hat\Delta_q \in [28, 67]$ bits pour les institutions bien pr\'epar\'ees
\item D\'ecroissance de $\Delta_q$ avec le niveau de bruit quantique
\item Corr\'elation positive entre formation PQC et $\Delta_q$
\end{enumerate}
\end{sloppypar}

% -----------------------------------------------------------------------
\subsection{Protocole de mesure de $\eta_q$~:
          coefficient de perte d'interpr\'etabilit\'e quantique}
\label{subsec:protocole_eta_q}
% -----------------------------------------------------------------------
\begin{sloppypar}

\subsubsection{Objectif et d\'efinition op\'erationnelle}
\label{subsubsec:objectif_eta_q}

\noindent\textbf{Objectif.} Mesurer empiriquement le coefficient de perte
d'interpr\'etabilit\'e quantique $\eta_q$, d\'efini
par~\cite{MINKA2026PQCRO}~:
\begin{equation}
I(V^q; \Lambda(\rho_\Sigma)|C)_\rho \;\geq\; I(V^q; \rho_\Sigma|C)_\rho - \eta_q \cdot H(C)
\label{eq:perte_interpretabilite_op}
\end{equation}
($\eta_q$ est le plus petit coefficient pour lequel cette minoration est
vraie~; dans l'autre sens, l'in\'egalit\'e serait satisfaite pour $\eta_q = 0$
par l'in\'egalit\'e de traitement des donn\'ees, ce qui viderait la d\'efinition de son contenu).

\noindent\textbf{Port\'ee.} Les signatures ML-DSA, SLH-DSA et Falcon produisent
des cha\^{\i}nes de bits classiques, pour lesquelles $\eta_q = 0$. Le protocole
ne s'applique qu'\`a un sc\'enario de stockage ou de v\'erification quantique,
qui reste hypoth\'etique \`a ce jour.

\noindent\textbf{D\'efinition op\'erationnelle.} $\eta_q$ quantifie la perte
d'information s\'emantique lorsqu'une preuve cryptographique traverse un canal
quantique (stockage, transmission, v\'erification avec ressources quantiques).

\noindent\textbf{Estimateur.}
\begin{equation}
\hat\eta_q = \frac{I(V^q; \rho_\Sigma|C)_\rho - I(V^q; \Lambda(\rho_\Sigma)|C)_\rho}{H(C)}
\label{eq:estimateur_eta_q}
\end{equation}

\subsubsection{Design exp\'erimental}
\label{subsubsec:design_eta_q}

\noindent\textbf{Principe.} Comparer l'extraction d'information depuis des
preuves avec et sans d\'egradation quantique simul\'ee.

\noindent\textbf{Mat\'eriel.}
\begin{itemize}
\item Simulateur de canaux quantiques (d\'epolarisation, d\'ephasage,
      amortissement d'amplitude)
\item Param\`etres de bruit calibr\'es sur les sp\'ecifications de mat\'eriel
      quantique actuel (IBM Quantum, IonQ, etc.)
\item Protocoles PQC avec et sans assistance quantique
\end{itemize}

\begin{table}[htbp]
\centering
\caption{Param\`etres de canaux quantiques pour la mesure de $\eta_q$}
\label{tab:val_canaux_quantiques}
\begin{tabular}{@{}lccc@{}}
\toprule
\textbf{Canal} & \textbf{Param\`etre} & \textbf{Plage test\'ee} & \textbf{$\eta_q$ attendu} \\
\midrule
D\'epolarisation & $p_{\mathrm{dep}}$ & $[0{,}001;\,0{,}1]$ & $0{,}02$--$0{,}20$ \\
D\'ephasage & $p_{\mathrm{deph}}$ & $[0{,}001;\,0{,}05]$ & $0{,}01$--$0{,}10$ \\
Amortissement & $\gamma_{\mathrm{AD}}$ & $[0{,}001;\,0{,}02]$ & $0{,}005$--$0{,}05$ \\
Combin\'e (NISQ typique) & --- & --- & $0{,}05$--$0{,}15$ \\
\bottomrule
\end{tabular}
\end{table}

\subsubsection{Proc\'edure}
\label{subsubsec:procedure_eta_q}

\noindent\textbf{Phase 1. Calibration.}
\'Etablir la ligne de base sans bruit quantique en mesurant
$I(V^q; \rho_\Sigma|C)_\rho$ pour chaque type de preuve.

\noindent\textbf{Phase 2. D\'egradation contr\^ol\'ee.}
Appliquer des canaux quantiques avec param\`etres croissants~:
\begin{equation}
\Lambda_{\mathrm{dep}}(\rho) = (1-p)\rho + \frac{p}{d}I_d
\label{eq:canal_depolarisant_op}
\end{equation}
Mesurer $I(V^q; \Lambda(\rho_\Sigma)|C)_\rho$ pour chaque niveau.

\noindent\textbf{Phase 3. Analyse de r\'egression.}
Ajuster le mod\`ele~:
\begin{equation}
\Delta I = I(V^q; \rho_\Sigma|C) - I(V^q; \Lambda(\rho_\Sigma)|C) = \eta_q \cdot H(C)
\end{equation}
pour extraire $\hat\eta_q$.

\subsubsection{M\'etriques compl\'ementaires}
\label{subsubsec:metriques_eta_q}

\begin{itemize}
\item \textbf{Fid\'elit\'e moyenne~:} $F_{\mathrm{avg}}(\Lambda) = \mathbb{E}[F(\Lambda(\rho), \rho)]$
\item \textbf{Distance de trace~:} $\|\Lambda(\rho_\Sigma) - \rho_\Sigma\|_1$
\item \textbf{Perte d'information par bit contextuel~:} $\Delta I / H(C)$
\end{itemize}

\subsubsection{Pr\'edictions v\'erifiables pour $\eta_q$}
\label{subsubsec:predictions_eta_q}

\begin{enumerate}
\item $\hat\eta_q \in [0{,}05;\,0{,}15]$ pour un support simul\'e aux taux d'erreur
      du mat\'eriel actuel (valeurs a priori, \`a pr\'e-enregistrer)
\item Ordre de grandeur de $\hat\eta_q\,H(C)$ compatible avec la borne de
      continuit\'e de l'information mutuelle conditionnelle (Alicki--Fannes--Winter),
      exprim\'ee en fonction de la distance de trace
      $\tfrac12\|\Lambda(\rho_\Sigma)-\rho_\Sigma\|_1 \leq \sqrt{1-F}$ (heuristique)
\item Relation croissante, approximativement lin\'eaire pour bruit faible, entre
      $p_{\mathrm{dep}}$ et $\hat\eta_q$
\item $\hat\eta_q = 0$, aux fluctuations d'\'echantillonnage pr\`es, lorsque la
      preuve est lue comme cha\^{\i}ne de bits (contr\^ole n\'egatif)
\end{enumerate}
\end{sloppypar}

% -----------------------------------------------------------------------
\subsection{Protocole combin\'e de validation quantique}
\label{subsec:protocole_combine}
% -----------------------------------------------------------------------
\begin{sloppypar}

\subsubsection{Objectif}
\label{subsubsec:objectif_combine}

Valider conjointement les param\`etres post-quantiques $(\Delta_q, \eta_q)$ et
comparer les pr\'edictions du Trilemme CRO post-quantique avec les mesures
empiriques sur les standards NIST PQC.

\subsubsection{Design int\'egr\'e}
\label{subsubsec:design_combine}

\begin{algorithm}[H]
\caption{Protocole combin\'e post-quantique}
\label{alg:protocole_combine}
\begin{algorithmic}[1]
\State \textbf{Entr\'ees~:} Institution $I$, standards PQC, canaux quantiques
\State \textbf{Sorties~:} $\hat\Delta_q$, $\hat\eta_q$, $\hat\beta_q$
\State \textbf{Partie A. Mesure de $\Delta_q$}
\State Mesurer extraction sur ML-DSA et SLH-DSA sans bruit
\State Calculer $\hat\Delta_q$ par comparaison avec gold standard PQC
\State \textbf{Partie B. Mesure de $\eta_q$}
\State Appliquer canaux quantiques \`a intensit\'e croissante
\State Mesurer d\'ecroissance de l'information mutuelle
\State Calculer $\hat\eta_q$ par r\'egression
\State \textbf{Partie C. Validation de la borne CRO$^q$}
\State V\'erifier les hypoth\`eses du th\'eor\`eme (dont $H(V|C) \geq \gamma$)
\State Si elles tiennent, confronter $\Conf^q \cdot \Relrec^q \cdot \Opp^q$ \`a la borne~\eqref{eq:trilemme_pq_raffine}
\State \textbf{Partie D. Validation crois\'ee}
\State Comparer pr\'ediction vs mesure pour les deux standards PQC
\State Calculer erreur relative et intervalles de confiance
\end{algorithmic}
\end{algorithm}

L'algorithme structure la validation post-quantique en quatre parties
ind\'ependantes, chacune mesurant ou v\'erifiant une composante distincte
de la borne CRO$^q$. L\`a o\`u les algorithmes pr\'ec\'edents isolaient
un param\`etre, celui-ci les articule puisqu'il produit les trois estimateurs
$(\hat\Delta_q, \hat\eta_q, \hat\beta_q)$, puis les confronte \`a la
borne th\'eorique et aux pr\'edictions sur les standards NIST PQC.

La phase d'\textbf{entr\'ees et sorties} (lignes~1--2) fixe le cadre.
Elle re\c{c}oit l'institution cible $I$, les standards PQC \`a tester et
les canaux quantiques disponibles. Elle d\'eclare trois sorties, \`a savoir
$\hat\Delta_q$ (capacit\'e d'extraction quantique), $\hat\eta_q$ (perte
d'interpr\'etabilit\'e) et $\hat\beta_q$ (r\'egularit\'e du protocole en
r\'egime quantique). Cette troisi\`eme sortie est nouvelle par rapport
aux algorithmes pr\'ec\'edents, car elle correspond au param\`etre $\beta$
de l'hypoth\`ese de r\'egularit\'e, transpos\'e au r\'egime quantique. Sa
mesure est n\'ecessaire pour v\'erifier que la borne post-quantique
s'applique bien aux protocoles test\'es.

La \textbf{partie~A} (lignes~3--5) mesure $\Delta_q$ en l'absence de
bruit. Elle \'evalue l'extraction institutionnelle sur les deux
sch\'emas de signature PQC (ML-DSA, SLH-DSA) dans leur
environnement nominal, sans d\'egradation quantique ajout\'ee. Le terme
$\hat\Delta_q$ est ensuite obtenu par comparaison avec un
\emph{gold standard} PQC, constitu\'e des verdicts de r\'ef\'erence \'etablis par
des experts en cryptographie post-quantique. Cette comparaison isole la
capacit\'e d'extraction proprement institutionnelle de toute
d\'egradation li\'ee au canal. La partie~A produit donc la valeur de
$\Delta_q$ qui serait observ\'ee sur un syst\`eme PQC id\'eal, sans
bruit.

La \textbf{partie~B} (lignes~6--9) mesure $\eta_q$ par d\'egradation
contr\^ol\'ee. Elle applique des canaux quantiques \`a intensit\'e
croissante, typiquement un canal d\'epolarisant dont le param\`etre
$p$ varie, puis observe la d\'ecroissance de l'information mutuelle
extraite. Cette d\'ecroissance, non lin\'eaire en $p$,
d\'epend de la structure du protocole et de la nature du canal. La
r\'egression sur les niveaux de bruit fournit $\hat\eta_q$ comme pente
de la relation entre perte d'information et bruit. La partie~B isole
donc l'effet du canal, ind\'ependamment de la capacit\'e
institutionnelle mesur\'ee en partie~A.

La \textbf{partie~C} (lignes~10--12) confronte les mesures \`a la borne
CRO$^q$. Elle v\'erifie d'abord que le cas \'etudi\'e satisfait les hypoth\`eses
du th\'eor\`eme, faute de quoi aucune comparaison n'est probante
(Remarque~\ref{rem:portee_cas_etude}). Si elles tiennent, le produit
$\Conf^q \cdot \Relrec^q \cdot \Opp^q$ est compar\'e \`a la forme raffin\'ee de la
borne (\'equation~\eqref{eq:trilemme_pq_raffine}). Un d\'epassement, hypoth\`eses
v\'erifi\'ees, r\'efute l'une d'elles selon les crit\`eres du
tableau~\ref{tab:criteres_refutation}~; il ne peut pas \^etre absorb\'e par un
r\'eajustement a posteriori de $\hat\Delta_q$.

La \textbf{partie~D} (lignes~13--15) proc\`ede \`a la validation
crois\'ee. Elle compare les pr\'edictions th\'eoriques aux mesures
empiriques pour chacun des deux standards PQC, puis calcule l'erreur
relative et les intervalles de confiance. Au-del\`a d'un contr\^ole de
coh\'erence, cette comparaison permet d'identifier les
standards pour lesquels le mod\`ele pr\'edit correctement les
compromis ainsi que ceux pour lesquels des ph\'enom\`enes non captur\'es par
le mod\`ele interviennent. Les intervalles de confiance quantifient
l'incertitude r\'esiduelle, qui est plus \'elev\'ee en r\'egime
post-quantique en raison de la variabilit\'e suppl\'ementaire introduite
par $\eta_q$.

L'algorithme se distingue des pr\'ec\'edents sur trois points. Il
articule quatre parties ind\'ependantes au lieu de mesurer un seul
param\`etre. Il produit un troisi\`eme estimateur $\hat\beta_q$. Enfin, il valide explicitement la borne
CRO$^q$ au lieu de se contenter de mesurer ses composantes. Cette articulation
fait du protocole combin\'e l'outil de validation complet du r\'egime
post-quantique.

\subsubsection{Crit\`eres de validation quantique}
\label{subsubsec:criteres_validation_quantique}

Les crit\`eres de r\'efutation de l'ensemble du programme, classique et
post-quantique, sont fix\'es avant toute collecte dans le
tableau~\ref{tab:criteres_refutation}. Les valeurs seuils sont indicatives et
seront arr\^et\'ees lors du pr\'e-enregistrement~; leur forme est
l'engagement m\'ethodologique.

\begin{table}[htbp]
\centering
\caption{Crit\`eres de r\'efutation pr\'e-enregistr\'es}
\label{tab:criteres_refutation}
\small
\begin{tabular}{@{}p{3.4cm}p{9.8cm}@{}}
\toprule
\textbf{Param\`etre} & \textbf{L'hypoth\`ese est r\'efut\'ee si\ldots} \\
\midrule
$\Delta$ (Hyp.~\ref{ass:extraction_bornee}) & la borne inf\'erieure de l'IC~95\,\% de $\hat{I}(V;\Sigma|M,C)$ d\'epasse la capacit\'e fix\'ee \emph{avant} l'exp\'erience pour la population \'etudi\'ee. \\
Admissibilit\'e (D\'ef.~\ref{def:interp_admissible}) & un test de permutation rejette $V \perp M \mid (\Sigma,C)$ au seuil $0{,}01$. \\
$\gamma$ (Hyp.~\ref{ass:verdicts_efficace}) & la borne sup\'erieure de l'IC~95\,\% de $\hat{H}(V|C)$ est inf\'erieure \`a $\alpha H(\mathcal{C}) - 1$. \\
$\eta_q$ (sc\'enario quantique) & l'IC~95\,\% de $\hat{\eta}_q$ exclut l'intervalle pr\'edit avant l'exp\'erience. \\
Borne (hypoth\`eses v\'erifi\'ees) & le produit mesur\'e d\'epasse la forme raffin\'ee de la borne avec une confiance de 95\,\%. \\
\bottomrule
\end{tabular}
\end{table}

\end{sloppypar}

\begin{boiteArbitre}
L'arbitre fait le point. Trois param\`etres classiques, trois param\`etres
quantiques, soit six protocoles qui visent tous \`a transformer des hypoth\`eses
en mesures. Il retient surtout leur sym\'etrie, puisque chaque quantit\'e
admise plus haut dispose d\'esormais d'une proc\'edure pour \^etre test\'ee. La
th\'eorie, d\'ej\`a coh\'erente, devient ainsi v\'erifiable.
\end{boiteArbitre}

% -----------------------------------------------------------------------
\subsection{Calendrier, ressources et statut de d\'eploiement}
\label{subsec:calendrier_ressources}
% -----------------------------------------------------------------------
\begin{sloppypar}

\subsubsection{Tableau de synth\`ese}
\label{subsubsec:tableau_synthese}

\begin{table}[htbp]
\centering
\caption{Programme exp\'erimental complet~: calendrier et ressources}
\label{tab:programme_experimental}
\begin{threeparttable}
\begin{tabular}{@{}lllcc@{}}
\toprule
\textbf{Protocole} & \textbf{Section} & \textbf{Dur\'ee} &
\textbf{Participants} & \textbf{Co\^ut estim\'e\tnote{$\star$}} \\
\midrule
\multicolumn{5}{c}{\textit{Param\`etres classiques (Contribution C2)}} \\
\midrule
Mesure de $\Delta$ & \S\ref{subsec:protocole_delta} & 6 mois & 120 institutionnels & 15\,000\,000 FCFA \\
Mesure de $H(\mathcal{C})$ & \S\ref{subsec:protocole_HC} & 3 mois & 20 experts & 5\,000\,000 FCFA \\
Mesure de $\Opp$ sur Helios & \S\ref{subsec:protocole_opp} & 4 mois & 10 trustees & 8\,000\,000 FCFA \\
\midrule
\multicolumn{5}{c}{\textit{Param\`etres quantiques (Contribution C3)}} \\
\midrule
Mesure de $\Delta_q$ & \S\ref{subsec:protocole_delta_q} & 6 mois & 80 professionnels PQC & 12\,000\,000 FCFA \\
Mesure de $\eta_q$ & \S\ref{subsec:protocole_eta_q} & 4 mois & 40 experts + simulation & 8\,000\,000 FCFA \\
Validation combin\'ee PQC & \S\ref{subsec:protocole_combine} & 3 mois & 40 validateurs & 6\,000\,000 FCFA \\
\midrule
\textbf{Total} & & \textbf{26 mois} & \textbf{310} &
\textbf{54\,000\,000 FCFA} \\
\bottomrule
\end{tabular}
\begin{tablenotes}
\footnotesize
\item[$\star$] Les co\^uts indiqu\'es constituent des ordres de grandeur
indicatifs destin\'es \`a fournir une vision globale du programme~;
ils feront l'objet d'une \'evaluation budg\'etaire affin\'ee lors de la
mise en \oe uvre effective de chaque protocole.
\end{tablenotes}
\end{threeparttable}
\end{table}

\subsubsection{Calendrier d\'etaill\'e}
\label{subsubsec:calendrier_detaille}

\begin{table}[htbp]
\centering
\caption{Calendrier du programme exp\'erimental}
\label{tab:calendrier_experimental}
\begin{tabular}{@{}lcccccc@{}}
\toprule
\textbf{Activit\'e} & \textbf{M1--M6} & \textbf{M7--M12} & \textbf{M13--M18} &
\textbf{M19--M24} & \textbf{M25--M30} & \textbf{M31--M36} \\
\midrule
Mesure $\Delta$ (\S\ref{subsec:protocole_delta}) & $\times\times$ & $\times\times$ & & & & \\
Mesure $H(\mathcal{C})$ (\S\ref{subsec:protocole_HC}) & $\times$ & $\times$ & & & & \\
Mesure $\Opp$ Helios (\S\ref{subsec:protocole_opp}) & & $\times\times$ & $\times$ & & & \\
Mesure $\Delta_q$ (\S\ref{subsec:protocole_delta_q}) & & & $\times\times$ & $\times\times$ & & \\
Mesure $\eta_q$ (\S\ref{subsec:protocole_eta_q}) & & & & $\times\times$ & $\times$ & \\
Validation PQC (\S\ref{subsec:protocole_combine}) & & & & & $\times\times$ & $\times$ \\
Analyse longitudinale & & & & & & $\times\times$ \\
\bottomrule
\end{tabular}
\end{table}

\subsubsection{Statut de d\'eploiement du programme}
\label{subsubsec:statut_deploiement}

Le programme d\'ecrit aux sections~\ref{subsec:protocole_delta}
\`a~\ref{subsec:calendrier_ressources} constitue une sp\'ecification
m\'ethodologique. Aucun de ses six
protocoles n'a \'et\'e d\'eploy\'e dans le cadre de cette th\`ese, qu'il s'agisse de la
campagne de mesure de $\Delta$ aupr\`es des 120 acteurs institutionnels
(section~\ref{subsec:protocole_delta}), de l'audit post-\'electoral
Helios \`a l'ENSP (section~\ref{subsec:protocole_opp}) ou des sessions
post-quantiques sur ML-DSA et SLH-DSA
(sections~\ref{subsec:protocole_delta_q}
et~\ref{subsec:protocole_eta_q}).

Ce choix d\'ecoule du positionnement \'epist\'emologique retenu au
chapitre~\ref{chap:cadre_formel}, selon lequel ce travail est d'abord th\'eorique. La
certification m\'ecanique en Lean4 garantit la correction de la d\'erivation \`a
partir des hypoth\`eses. La v\'erit\'e des hypoth\`eses elles-m\^emes rel\`eve de la
validation comportementale, que seules des mesures peuvent \'etablir. Trois pr\'ealables
non techniques conditionnent aussi le d\'eploiement, \`a savoir l'accord des
juridictions et r\'egulateurs approch\'es, l'approbation d'un comit\'e
d'\'ethique de la recherche ainsi que le financement de l'ordre de grandeur
report\'e au tableau~\ref{tab:programme_experimental}, pour vingt-six mois de
travail cumul\'es r\'epartis sur un calendrier de trente-six mois
(tableau~\ref{tab:calendrier_experimental}).

Il en r\'esulte, pour la lecture des sections suivantes, que les
d\'erivations et les estimations qui y figurent sont de
nature r\'etrospective et analytique. Cela vaut pour les valeurs
d'opposabilit\'e d\'eriv\'ees \`a la
section~\ref{sec:derivation_opposabilite} comme pour les param\`etres
institutionnels estim\'es \`a la
section~\ref{sec:validation_retrospective}, puisque les unes et les autres
reposent sur les sp\'ecifications publi\'ees des protocoles et sur des
contextes institutionnels mod\'elis\'es, sans mesure aupr\`es de juges ou
de r\'egulateurs r\'eels. Elles \'etablissent que les param\`etres
institutionnels admettent des valeurs plausibles et mutuellement
coh\'erentes. Seule une mesure ind\'ependante dira si ces valeurs sont les
bonnes. L'audit
post-\'electoral Helios, le plus l\'eger en ressources et le seul
enti\`erement internalisable \`a l'ENSP, constitue le point d'entr\'ee
naturel du programme~; il est repris \`a ce titre dans les perspectives
du chapitre~\ref{chap:resultats_discussions}.

\end{sloppypar}

\begin{boiteArbitre}
L'arbitre referme le dossier des protocoles. Il n'a rien ex\'ecut\'e, rien
mesur\'e, rien valid\'e. Il tient une promesse de mesure rigoureuse, chiffr\'ee et falsifiable, qui reste une promesse.

Il a compris que ces protocoles \'etablissent le caract\`ere \emph{testable}
du cadre, ce qui est moins que sa v\'erit\'e sans \^etre peu. Une th\'eorie qu'on
ne peut pas tester reste sans prise sur le r\'eel, m\^eme si elle est belle.

Il rel\`eve la transparence du dispositif, dont les pr\'edictions
sont num\'eriques, les crit\`eres d'invalidation explicites et les
limites nomm\'ees. L'audit Helios lui-m\^eme est pr\'esent\'e comme le
point d'entr\'ee le plus modeste, qui est aussi le seul internalisable \`a l'ENSP.

Il ne sait pas encore si les chiffres tiennent, mais il sait d\'esormais
\emph{comment} il le saura, ce qui correspond exactement \`a ce qu'il attendait d'un
protocole exp\'erimental.

Les d\'erivations peuvent commencer.
\end{boiteArbitre}